% FORMALIZACION_REUNIDA de la proyección del calendario y la representación
% de doce posiciones ya desarrolladas en la arquitectura TPK y el enlace Witt.
% Formalización explícita del cotejo de Gram y del transporte de marco.
\subsubsection{Proyección del calendario y representación de sus doce posiciones}
\label{ii:dodeca:calendario}
El calendario de la sección~\ref{ii:tpk-arquitectura} proporciona una
coordenada observable del estado completo. Con sus notaciones, escribimos
\(r_0(\theta)=r(\theta)-1\in\{0,\ldots,8\}\) y
\(b=\beta(\theta)\in\mathbb Z/12\mathbb Z\). La proyección de reloj es
\[
 Q_{\rm clk}(x)=\bigl(\beta(\theta(x)),r_0(\theta(x))\bigr).
\]
Para una ruta admisible de \(n\) actualizaciones, el transporte inducido es
\begin{equation}
 k(r_0,n)=\left\lfloor\frac{r_0+n}{9}\right\rfloor,
 \qquad
 T_n(b,r_0)=\bigl(b+k(r_0,n)\bmod12,\ (r_0+n)\bmod9\bigr).
 \label{ii:dodeca:transporte-reloj}
\end{equation}
Aquí \(n\in\mathbb N_0\) cuenta actualizaciones efectivas; el cociente
entero \(k\) registra los retornos nonádicos atravesados.

\begin{proposition}[Compatibilidad del calendario con la composición de rutas]
\label{ii:dodeca:composicion}
La proyección \(Q_{\rm clk}\), junto con la longitud de cada ruta, define
un funtor de la categoría de rutas admisibles TPK a la categoría de acción
del calendario~\eqref{ii:dodeca:transporte-reloj}. Se cumplen
\(T_{m+n}=T_mT_n\), \(T_0=I\),
\(T_9(b,r_0)=(b+1,r_0)\) y \(T_{108}(b,r_0)=(b,r_0)\).
\end{proposition}
\begin{proof}
Poniendo \(r'_0=(r_0+n)\bmod9\), la división euclidiana da
\[
 k(r_0,n+m)=k(r_0,n)+k(r'_0,m).
\]
Esta identidad conserva simultáneamente residuo y cociente al concatenar
las rutas. La actualización de \(\theta\) establecida en la
sección~\ref{ii:tpk-arquitectura} induce exactamente
\eqref{ii:dodeca:transporte-reloj}; las identidades restantes se obtienen
por sustitución. La categoría de llegada tiene los pares \((b,r_0)\) como
objetos y las longitudes admisibles como flechas, con suma como composición.
\end{proof}

El índice \(p=b+1\in\{1,\ldots,12\}\), tomando el representante
\(0\le b<12\), indica la posición dodecafásica. Una vuelta de
\(108\) actualizaciones recorre sus doce posiciones. La periodicidad
demostrada corresponde a esta proyección: ruta, acarreo acumulado y
memoria permanecen en el estado fuente y continúan su actualización.

\paragraph{Realización de incidencia a partir del proyector angular.}
Los representantes \(r_pC_5\) y la representación \(\rho\) de la
sección~\ref{sec:ii-enlace-angular-witt} fijan la coordenatización de las doce
posiciones. El proyector ya definido por
\eqref{eq:ii-witt-projector-character} produce
\begin{equation}
 \begin{gathered}
 \mathcal H_3=P_3^{\mathrm{ico}}\mathbb R^{12},\qquad v_p=2P_3^{\mathrm{ico}}e_p,
 \\[-2pt]
 (P_3^{\mathrm{ico}})_{pq}=\frac1{20}
 \sum_{\substack{g\in A_5\\gr_qC_5=r_pC_5}}\chi_3(g^{-1}).
 \end{gathered}
 \label{ii:dodeca:vertices-proyector}
\end{equation}
La última suma contiene cinco términos por entrada y utiliza los
representantes y las dos clases de ciclos de orden cinco allí fijados.

\begin{proposition}[Realización icosaédrica de las posiciones dodecafásicas]
\label{ii:dodeca:realizacion}
Los vectores \(v_p\) son doce vectores unitarios distintos de
\(\mathcal H_3\). Su incidencia es
\(p\sim q\Longleftrightarrow\langle v_p,v_q\rangle=1/\sqrt5\).
La aplicación \(p\mapsto v_p\) es \(A_5\)-equivariante y su imagen
es isométrica a \(\Omega_{12}^{\rm ico}\).
\end{proposition}
\begin{proof}
La transitividad y \(\operatorname{Tr}P_3^{\mathrm{ico}}=3\) dan
\((P_3^{\mathrm{ico}})_{pp}=1/4\); por idempotencia,
\(\langle v_p,v_q\rangle=4(P_3^{\mathrm{ico}})_{pq}\) y \(\|v_p\|=1\).
La evaluación de los cinco términos de
\eqref{ii:dodeca:vertices-proyector} da los pares antipodales
\[
 (1,4),\ (2,3),\ (5,8),\ (6,7),\ (9,12),\ (10,11).
\]
Para los representantes \((1,2,5,6,9,10)\), su matriz de Gram es
\begin{equation}
 I_6+\frac D{\sqrt5},\qquad
 D=\begin{pmatrix}
 0&-1&-1&1&-1&1\\
 -1&0&1&1&-1&-1\\
 -1&1&0&1&1&1\\
 1&1&1&0&-1&1\\
 -1&-1&1&-1&0&1\\
 1&-1&1&1&1&0
 \end{pmatrix},\qquad D^2=5I_6.
 \label{ii:dodeca:gram}
\end{equation}
Los restantes productos escalares se obtienen mediante antipodía.
Así, aparte de un vértice y su opuesto, los productos son
\(\pm1/\sqrt5\), y cada vértice tiene cinco vecinos. La tabla siguiente
explicita el cotejo con la coordenatización firmada de seis posiciones
\(\mathcal U=(\infty,0,1,2,3,4)\):
\begin{equation}
\begin{array}{c|cc@{\qquad}c|cc}
p&u(p)&\sigma(p)&p&u(p)&\sigma(p)\\ \hline
1&\infty&+1&7&3&-1\\
2&0&-1&8&1&+1\\
3&0&+1&9&2&-1\\
4&\infty&-1&10&4&+1\\
5&1&-1&11&4&-1\\
6&3&+1&12&2&+1
\end{array}
\label{ii:dodeca:carta-firmada}
\end{equation}
En efecto, para \(C=\operatorname{bal}(A_W)\), con \(A_W\) dado en
\eqref{exc:aw}, las entradas de \(D\) son
\(D_{ab}=\sigma(p_a)\sigma(p_b)C_{u(p_a),u(p_b)}\).
Por consiguiente, la tabla y~\eqref{ii:dodeca:gram} identifican
isométricamente \(v_p\) con
\(\sigma(p)\sqrt2 E_Ce_{u(p)}\), donde
\(E_C=(I_6+C/\sqrt5)/2\). Ambos sistemas generan espacios de dimensión
tres y tienen idéntica matriz de Gram, lo que prueba que la aplicación
lineal inducida está bien definida y es isométrica. La realización
explícita de~\eqref{ii:ico:proyector-seis}, demostrada a continuación,
la identifica con~\(\Omega_{12}^{\rm ico}\) y su incidencia.
Finalmente, \(P_3^{\mathrm{ico}}\rho(g)=\rho(g)P_3^{\mathrm{ico}}\) implica
\(v_{g\cdot p}=\rho(g)v_p\). La biyección de las doce posiciones da
\(\operatorname{Stab}_{A_5}(v_p)=r_pC_5r_p^{-1}\).
\end{proof}

Se obtiene así la composición representacional
\begin{equation}
 x\longmapsto Q_{\rm clk}(x)\longmapsto
 p=\beta(\theta(x))+1\longmapsto v_p.
 \label{ii:dodeca:composicion-representacional}
\end{equation}
La expresión firmada de esta composición es
\(r_\pm^{\rm rep}(x)=(u(p),\sigma(p))\), determinada por la
tabla~\eqref{ii:dodeca:carta-firmada}. La tabla es una elección explícita
de alineamiento entre coordenatizaciones; las acciones de \(A_5\) proporcionan sus
alineamientos equivalentes. El registro completo se conserva en el
dominio de~\eqref{ii:dodeca:composicion-representacional}.

\paragraph{Enlace con el portador tensorial de cinco componentes.}
En \(\mathcal H_3\), formamos
\(Q_p=v_pv_p^*-I_{\mathcal H_3}/3\). Sea \(J=\mathbf1\mathbf1^*\)
y sea \(\mathsf A\) la matriz de la permutación antipodal anterior.
El producto de Frobenius y~\eqref{ii:dodeca:gram} dan
\begin{equation}
 \begin{aligned}
 \operatorname{Gram}(Q_p)
 &=(4P_3^{\mathrm{ico}})\circ(4P_3^{\mathrm{ico}})-\frac13J
 =\frac45(I+\mathsf A)-\frac2{15}J,
 \\
 P^{\rm ico}_5&=\frac58\operatorname{Gram}(Q_p)
 =\frac12(I+\mathsf A)-\frac1{12}J.
 \end{aligned}
 \label{ii:dodeca:puente-tensorial}
\end{equation}
El símbolo \(\circ\) denota el producto entrada a entrada. La última
expresión es el proyector ortogonal sobre las coordenadas pares bajo
antipodía y ortogonales a \(\mathbf1\), de dimensión cinco. Su carácter
es \((5,1,-1,0,0)\): sobre los seis pares antipodales, los números de
puntos fijos son \((6,2,0,1,1)\), y se retira la componente constante.
Por ello coincide con el proyector central de
\eqref{ii:pent:eq:gravity-p5}. Esta identidad une la construcción
tridimensional con \(\Gamma_0\) y \(\Gamma_5\), definidos en
\eqref{ii:pent:eq:gravity-Gamma0} y~\eqref{ii:pent:eq:gravity-gamma5},
mediante la operación cuadrática explícita \(v_p\mapsto Q_p\).

\subsubsection{Transporte del marco representacional}
\label{ii:dodeca:marco}
Sea \(s=(1\,2\,\cdots\,12)\) y sea \(S_ce_p=e_{s(p)}\).
Un retorno nonádico del calendario lleva \(p\) a \(s(p)\). El transporte
conjunto de marco, representación y proyector se expresa por
\begin{equation}
 \begin{gathered}
 P_{3,k}^{\mathrm{ico}}=S_c^kP_3^{\mathrm{ico}}S_c^{-k},\qquad
 \rho^{(k)}(g)=S_c^k\rho(g)S_c^{-k},\\
 v^{(k)}_{s^k(p)}=2P_{3,k}^{\mathrm{ico}}e_{s^k(p)}=S_c^kv_p.
 \end{gathered}
 \label{ii:dodeca:marco-movil}
\end{equation}
La última igualdad se prueba por sustitución; la ortogonalidad de
\(S_c\) conserva todos los productos escalares y transporta la
incidencia. Para una ruta de longitud \(n\), se toma
\(k=k(r_0,n)\). La identidad de acarreo de la
proposición~\ref{ii:dodeca:composicion} prueba la compatibilidad de
estos transportes con la concatenación. Aunque \(S_c^{12}=I\),
las doce vueltas actualizan la memoria de la ruta fuente.

La igualdad~\eqref{ii:dodeca:marco-movil} establece covariancia entre
marcos. La acción activa sobre una incidencia fija tiene un contenido
adicional: requiere identificar el transporte del estado completo y su
acción en esa coordenatización. En particular, \(s\) tiene orden doce, mientras
que los órdenes de los elementos de \(A_5\) son \(1,2,3,5\).
Asimismo, el estabilizador de un vértice tiene orden cinco, y el de su
par antipodal tiene orden diez; la involución de la coordenatización firmada
intercambia los dos vértices. Estas distinciones conservan el significado
de orientación, hoja y retorno en las propiedades (i)--(iv): la
realización representacional demostrada y el descenso de la ruta
excepcional completa se relacionan mediante esas propiedades, con sus
dominios y fibras declarados.
